A la fecha, la \emph{World Wide Web (WWW)} contiene más de 5.470 millones de páginas Web indexadas\cite{www2020}, en su gran mayoría, conteniendo información escrita en lenguaje natural. Si tuviéramos que imprimir todo este contenido en papel, y siendo muy conservadores,\footnote{Suponiendo, ingenuamente, que cada página web requiere, en promedio, de 30 hojas para imprimirse\cite{harwood2015much}.} requeriríamos al menos de 164.100 millones de páginas A4 para lograrlo.
En este contexto, es beneficioso contar con sistemas computacionales que sean capaces de procesar, adquirir, extraer información útil y realizar tareas valiosas, a partir del texto en crudo.
El diseño y desarrollo de dichos sistemas es uno de los problemas más desafiantes dentro de la \ac{IA}, ya que estos sistemas necesitan entender, en cierta medida, el lenguaje ambiguo, desordenado, y poco estructurado usado por nosotros, los humanos.
Tal es la relevancia de este desafiante problema que su origen puede rastrearse incluso hasta los orígenes de la \ac{IA} como tal, ya que, cuando Alan Turing propone su famoso ``test de Turing''\cite{turing2009computing}\footnote{Originalmente llamado ``The Imitation Game''.} como definición pragmática de inteligencia, lo hace en relación a la capacidad de los sistemas para procesar y comprender el lenguaje natural de su interlocutor.
En este capítulo, nos restringiremos a examinar una tarea específica dentro de esta problemática, denominada como \emph{clasificación de texto}, que, si bien tal vez pueda parecer acotada, en la práctica tiene una amplia aplicabilidad.



\section{Introducción}

En la \autoref{sec:clasificacion} del capítulo anterior se introdujo el concepto de \emph{clasificación} dentro del marco del aprendizaje automático supervisado. Aquí nos centraremos en una tarea de clasificación específica, conocida como clasificación de texto:\footnote{También conocida como \emph{categorización de texto}.} \emph{dado un texto de entrada, decidir a cuál de un conjunto predefinido de clases pertenece}.
Algunos ejemplos de aplicaciones reales de clasificación de texto serían el análisis de sentimientos en reseñas de películas o productos (clasificar una reseña como positiva o negativa), filtrado de spam (clasificar un e-mail como ``spam'' o ``no spam''), identificación automática de idiomas, categorización temática (clasificar documentos de acuerdo a su temática), así como también, la tarea abordada en este trabajo doctoral, esto es, la detección anticipada de riesgos, tales como depresión, anorexia o conductas autodestructivas, en base al contenido generado por el usuario en redes sociales.

En el capítulo anterior también vimos que los modelos de clasificación trabajan con instancias de entrada representadas por vectores de $m$ atributos, $\pmb{x}=\langle x_1,x_2,\dots,x_m\rangle$.
Esto implica que, en principio, el texto no puede ser interpretado directamente por los modelos, no sin antes ser procesado y transformado adecuadamente. Para llevar esto a cabo, normalmente se sigue un proceso que puede ser dividido en, al menos, dos etapas principales: una etapa encargada de preprocesar el texto y otra relacionada a la transformación del mismo a una representación vectorial propiamente dicha. Así, este capítulo se organiza de la siguiente manera: en la siguiente sección, \autoref{sec:preprocesado}, se describe brevemente esta primer etapa de preprocesamiento, mientras que la segunda es abordada en la \autoref{sec:representacion}, en donde se examinan distintos enfoques para representar el texto como un vector de atributos. Finalmente, la \autoref{sec:EC} introduce formalmente la problemática abordada en este trabajo de investigación, enmarcándolo como un escenario particular de clasificación de texto.
Por este motivo, los lectores ya familiarizados con la clasificación de textos, pero no con la problemática abordada, pueden continuar su lectura directamente a partir de la última sección.

\section{Preprocesamiento}
\label{sec:preprocesado}

% En clasificación de texto, a cada instancia de entrada se la denomina típicamente como documento de entrada.
El preprocesamiento consiste básicamente en la preparación del texto para su posterior procesamiento.
De forma general, preprocesar el texto simplemente significa llevarlo a una forma ``estándar'' que sea predecible y analizable, lo que ``estándar'' realmente signifique dependerá de  la tarea abordada y, muchas veces, también del clasificador utilizado.
La elección del preprocesamiento adecuado puede mejorar significativamente el desempeño de la clasificación\cite{uysal2014impact,hacohen2020influence}.
Sin embargo, el preprocesamiento de texto no es directamente transferible ni de una tarea a otra, ni de un modelo a otro ---puede ocurrir que el preprocesamiento ideal para una tarea, o para un modelo, se convierta en la peor opción para otro.

Existen diferentes tipos de técnicas de preprocesamiento de texto que pueden aplicarse en simultáneo sobre el texto de entrada. Estas técnicas pueden variar desde lo simple, como ser una sencilla conversión del texto a letras minúsculas, hasta lo elaborado, como ser la corrección automática de errores ortográficos o el enriquecimiento del texto usando etiquetadores gramaticales.
Sin embargo, ya que el preprocesamiento consta típicamente de tres etapas bien definidas, las técnicas estarán asociadas a cada una de estas etapas del siguiente modo: 

\begin{itemize}
    \item\textbf{Eliminación de ruido:} como su nombre lo indica, esta etapa se encarga de la eliminación del ruido en el texto de entrada con el fin de obtener el texto crudo. De entre todas, esta etapa es la que más depende del problema específico que está siendo abordado, ya que dependerá del dominio en el que se esté trabajando y el formato del texto de entrada, lo que implique ruido o no.
    Ejemplos de eliminación de ruido podrían ser la eliminación de palabras clave específicas al dominio, como ser ``RT'' para retweet, la eliminación de metadatos, tales como encabezados específicos de ciertos archivos, o la eliminación de información de formato o presentación, como las etiquetas HTML/XML.
    
    \item\textbf{Tokenización:} en es
    ta etapa se divide el texto en partes más pequeñas o ``tokens''. Por ejemplo, trozos grandes de texto pueden tokenizarse en oraciones y las oraciones tokenizarse en palabras.
    %En ocasiones, a la tokenización también se la suele conocer como ``segmentación de texto'' o ``análisis léxico''.
    Muy a menudo se suele reservar el término ``tokenización'' exclusivamente para la división del texto en palabras individuales, y el término ``segmentación'' para la división en bloques más grandes, como ser párrafos u oraciones.
    En esta etapa se pueden utilizar distintas técnicas, tanto para identificar los límites de los distintos segmentos como para determinar qué hacer una vez se los identifica.
    En cualquier caso, una vez que el texto ha sido tokenizado adecuadamente, generalmente, se procede con el resto del preprocesamiento.\footnote{En ciertos idiomas, por ejemplo, idiomas aglutinantes como el alemán o el japonés, esta etapa también tiene que ver con separar palabras compuestas en sus monemas independientes.}
    
    \item\textbf{Normalización:} esta etapa se refiere, generalmente, a la aplicación de una serie de tareas con el fin de dejar a todo el texto en ``igualdad de condiciones''.
    De este modo, la normalización permite que el procesamiento de las entradas proceda de manera uniforme, independientemente de cómo eran dichas entradas originalmente.
    Esta etapa, normalmente, involucra alguna(s) de las siguientes tareas: convertir todo el texto a minúsculas (o mayúsculas), eliminar las signos de puntuación, eliminar o convertir los números a su equivalente en palabras, eliminar las \emph{palabras vacías},\footnote{En el campo del \ac{PLN}, se utiliza el término ``palabra vacía'', en inglés ``stop word'', para referirse a una palabra que carece de semántica en relación al problema abordado. A excepción de problemas muy específicos, las palabras vacías se corresponden con las ``palabras funcionales'' en lingüística, \emph{i.e.} palabras que desempeñan un papel dentro del lenguaje en sí mismo sin referir a la realidad extralingüística. Por este motivo, en la práctica es habitual la utilización del término ``palabra vacía'' para referirse a palabras tales como artículos, preposiciones o conjunciones.} eliminar palabras de muy baja frecuencia, enriquecer/aumentar el texto con información valiosa para el problema abordado y reemplazar palabras por su raíz.
    Siendo \emph{stemming} y \emph{lemmatization} las dos técnicas más utilizadas para llevar a cabo esta última tarea.
    El primero construye la raíz simplemente eliminando los afijos (típicamente los sufijos) de las palabras,\footnote{Por ejemplo, ``pensaba'' se reemplazaría por ``pens'' tras la eliminación del sufijo ``aba''.} mientras que el segundo, más elaborado, usa el lema de la palabra como su raíz.\footnote{Por ejemplo, las tres palabras, ``pensaba'', ``pensábamos'' y ``pensando'' se reemplazarían por el lema, ``pensar''.}
\end{itemize}


En la práctica, dependiendo de la tarea, cualquiera de estas etapas podrá o no aplicarse. De hecho, estas etapas no necesariamente representan un proceso lineal cuyos pasos deben aplicarse exclusivamente en el orden dado.

\section{Representación}
\label{sec:representacion}

Como se mencionó al comienzo, típicamente, el texto no puede ser interpretado directamente por un clasificador, por lo que una vez preprocesado, el siguiente paso a definir es cómo se lo representará para poder ser efectivamente ingresado al modelo de clasificación.
Esto es, se debe definir qué procedimiento se utilizará para mapear el texto de un documento, $d$, a una representación vectorial de su contenido, $d=\langle w_1,\dots,w_m\rangle$.
La elección de una representación adecuada suele ser uno de los aspectos más imperantes, sino el más, de la creación de un buen clasificador.
A su vez, los distintos tipos de representaciones dependen principalmente de lo que se consideren como las unidades significativas del texto y la semántica de cada una de ellas~\cite{turney2010frequency}.
Así, de acuerdo a los objetivos de esta tesis, nos limitaremos a identificar y describir dos tipos de representaciones fundamentales: representaciones dadas por léxico y representaciones dadas por contexto.
En la primera, el texto es modelado como una secuencia de unidades léxicas, mientras que la segunda intenta capturar el significado de las palabras en base al contexto en el que éstas aparecen. El resto de esta sección está destinada a examinar brevemente estos dos grupos de representaciones.


\subsection{Representaciones dadas por Léxico}


Según la Real Academia Española, se define a léxico como el conjunto de palabras que conforman un determinado vocabulario.
En este contexto, las representaciones dadas por léxico hacen referencia a las representaciones que se basan en la construcción de un vocabulario utilizando los documentos del conjunto de entrenamiento.
Este vocabulario no necesariamente estará formado por palabras, sino que, de forma general, estará formado por cualquier tipo de unidades léxicas, tales como trozos de palabras, palabras individuales o, incluso, cadenas de palabras.
En la clasificación de texto, a estas unidades léxicas se las conoce como \textbf{términos}, por lo que es habitual denominar a este vocabulario como ``vocabulario de términos''.

De esta forma, para construir la representación de un documento $d$, primero se procesa el texto como una secuencia de términos\footnote{Este secuencia de términos normalmente se crea dividiéndolo durante la etapa de tokenización del preprocesamiento.} y luego se lo convierte en un vector $d=\langle w_1,\dots,w_m\rangle$, donde $m$ es el número total de términos en el vocabulario, y $w_i$ es el valor, denominado como ``peso'', asociado al $i$-ésimo término del vocabulario.
Este peso representa, hablando vagamente, cuánto contribuye cada término en particular a la semántica del documento $d$.
Por lo que, distintas representaciones particulares surgirán de las diferentes formas de entender qué es un término y de las diferentes formas de calcular el pesado de los mismos. En las siguientes dos subsecciones, tituladas ``Tipos de Términos'' y ``Pesado de Términos'', se examinan respectivamente ambos aspectos.

\subsubsection{Tipos de Términos}
\label{sec:tipo-termino}

A la hora de representar documentos, sin duda, la elección más ampliamente utilizada es identificar a los términos con las palabras individuales. A esta representación se la conoce con el nombre particular de \textbf{bolsa de palabras} o simplemente como \emph{BoW}, del inglés \emph{Bag of Words}. 
Este nombre surge de que, bajo esta suposición, se puede pensar a cada documento como una bolsa en la que se han puesto todas sus palabras, una a una, y, luego, se utilizará esta última para construir la representación, en lugar del documento original.
A pesar de su simpleza, la representación de bolsa de palabras es uno de los  enfoques más efectivos \cite{nguyen2013old,stamatatos2009survey,sebastiani2002machine}, y en la práctica es muy común encontrar que representaciones más sofisticadas no produzcan una mejora significativa de rendimiento \cite{apte1994automated,dumais1998inductive,lewis1992evaluation}.
Sin embargo, a pesar de su amplio uso y efectividad, existen algunos inconvenientes que pueden provocar que su aplicabilidad no sea adecuada en algunos escenarios.
Por ejemplo, al utilizar esta representación, la noción de orden entre las palabras se pierde y, por lo tanto, el modelo clasificará de la misma manera a cualquier permutación posible del texto original y será incapaz de capturar secuencias de palabras importantes ---lo que para algunas tareas sí puede ser relevante. Además, en escenarios donde hay grandes vocabularios, pero pocos documentos de entrenamiento y clases desequilibradas, esta representación puede contribuir a favorecer a las clases mayoritarias por sobre las demás\cite{stamatatos2008author}.

Una de las formas más habituales para intentar mitigar algunas de las limitaciones de la representación de bolsa de palabras es identificar a los términos con secuencias de palabras, en lugar de con palabras individuales. De esta forma, la representación mantiene, en cada término, alguna noción local del orden de las palabras del documento original.
Dado que una secuencia de símbolos de longitud $n$ se llama como \emph{n-grama}, del griego \textgreek{γραμμα} (\emph{gramma}) letra o escrito, a esta representación se la conoce con el nombre de \textbf{representación de \emph{n-grama} de palabras}, siendo casos especiales ``unigrama'' para \emph{1-grama},\footnote{De aquí que a la representación de bolsa de palabras también se la pueda llamar como representación de unigrama de palabras.} ``bigrama'' para \emph{2-grama} y ``trigrama'' para \emph{3-grama}.
Si bien esta representación puede tener cualidades semánticas superiores con respecto a la representación de bolsa de palabras, también sufre de cualidades estadísticas inferiores, ya que, conforme mayor sea el orden del \emph{n-grama}, menor será la frecuencia con la que aparezcan los términos (\emph{i.e.} los \emph{n-gramas de palabras}) en los documentos de entrenamiento\cite{lewis1992evaluation}.\footnote{La probabilidad con la que aparece una secuencia de palabras nunca será mayor que la probabilidad individual de cada palabra.}

Finalmente, en última instancia, un texto escrito no es más que una secuencia de caracteres, por lo que tiene sentido representarlo como tal. Así, otra elección, también muy habitual, es identificar a los términos con secuencias de caracteres, en cuyo caso la representación es denominada como \textbf{representación de \emph{n-gramas} de caracteres}.
Por medio de esta representación, los clasificadores son capaces de capturar no sólo contenido \emph{per se}, sino también información estilística y contextual\cite{stamatatos2009survey}.
Además, también suelen ser más resistentes a los errores ortográficos que las representaciones basadas en palabras completas. Por ejemplo, las palabras ``Paris'' y ``Pariz'' serán consideradas como dos términos completamente distintos por estos últimos, mientras que utilizando una representación de, digamos, trigrama de caracteres, ambas palabras aportarían, en gran medida, el mismo valor a la clasificación final, ya que ambas serán vistas como los términos ``Par''  y ``ari'', más una discrepancia en los términos finales (``ris'' y ``riz'') dada por el error ortográfico.
Sin embargo, a pesar de estas ventajas, también existen algunas desventajas con respecto a las representaciones basadas en palabras completas, como ser la alta dimensionalidad en la representación vectorial final.

\

Antes de finalizar esta subsección, es adecuado resaltar nuevamente que, en las tareas de clasificación de textos, cada problema específico tendrá sus propias particularidades, por lo que podrá ocurrir que la representación ideal para una tarea no sea conveniente para otra.
Por ejemplo, mientras que en tareas de clasificación de tópicos los términos más valiosos suelen ser las palabras temáticas, capturadas fácilmente por representaciones de bolsa de palabras, en tareas de atribución de autoría o detección de plagio, lo suelen ser los n-gramas de caracteres capturando información estilística, ya que el objetivo principal es modelar el estilo de escritura de un autor\cite{stamatatos2009survey}.
 
\subsubsection{Pesado de Términos}

Hasta aquí se han visto los distintos tipos de términos que se pueden utilizar para representar a los documentos, en esta subsección examinaremos brevemente aspectos relacionados a cómo pesarlos, \emph{i.e.} vinculados a cómo calcular los valores $w_i$ asociados a cada término, en el vector que representa al documento, $d=\langle w_1,\dots,w_m\rangle$.

En la práctica, por diversos motivos, tales como herencia por legado del campo de la \emph{Recuperación de Información},\footnote{En las ciencias computacionales, es el campo que históricamente se encarga de la búsqueda y recuperación de documentos que son relevantes en relación a la información necesitada por un usuario dado.} conveniencia matemática y  eficiencia computacional, el cálculo de los pesos, $w_i$, de cada término se calcula en primera instancia construyendo una matriz, denominada como \textbf{matriz documento-término}, que agrupa los vectores de todos los documentos del conjunto de entrenamiento, uno encima del otro.
De esta manera, si se tienen $n$ documentos de entrenamiento y el tamaño del vocabulario es de $m$ términos, la matriz resultante será de $n\times m$, ya que las filas se corresponden a los documentos y las columnas a los términos del vocabulario, como se ilustra en la siguiente figura:

\begin{center}
\begin{tabular}{c|c|c|c|c|c|}
\multicolumn{1}{c}{} & \multicolumn{1}{c}{$t_1$} & \multicolumn{1}{c}{$t_2$} & \multicolumn{1}{c}{$\dots$} & \multicolumn{1}{c}{$t_{m-1}$} & \multicolumn{1}{c}{$t_m$} \\
\cline{2-6}
$d_1$ & & & & & \\
\cline{2-6}
$d_2$ & & & & & \\
\cline{2-6}
$\vdots$ & & & & & \\
\cline{2-6}
$d_{n-1}$ & & & & & \\
\cline{2-6}
$d_n$ & & & & & \\
\cline{2-6}
% \cline{2-6}
% $d_1$ & $w_{1,1}$ & $w_{1,2}$ & $\dots$ & $w_{1,m-1}$ & $w_{1,m}$\\
% \cline{2-6}
% $d_2$ & $w_{2,1}$ & $w_{2,2}$ & $\dots$ &  $w_{2,m-1}$ & $w_{2,m}$\\
% \cline{2-6}
% $\vdots$ & $\vdots$ & $\vdots$ & $\vdots$ & $\vdots$ & $\vdots$ \\
% \cline{2-6}
% $d_{n-1}$ & $w_{n-1,1}$ & $w_{n-1,2}$ & $\dots$ & $w_{n-1,m-1}$ & $w_{n-1,m}$\\
% \cline{2-6}
% $d_n$ & $w_{n,1}$ & $w_{n,2}$ & $\dots$ & $w_{n,m-1}$ & $w_{n,m}$\\
% \cline{2-6}
\end{tabular}
\end{center}

Donde $d_i$ y $t_j$ denotan los índices asociados, respectivamente, al $i$-ésimo documento y al $j$-ésimo término.
Esta matriz permite visualizar claramente la relación entre los términos y los documentos, ya que utilizando sus índices se puede obtener el valor $w_{i,j}$ para cualquier par documento-término.
% Naturalmente, la mayoría de los documentos no contendrán a la mayoría de los términos, por lo que esta matriz no será densa, sino dispersa.\footnote{Es decir, la mayor parte de sus elementos serán cero.}

Luego, distintos tipos de pesados de términos surgen de distintas maneras de completar esta matriz documento-término.
Así, una primera opción es asignarle un 1, o un 0, a cada posición, dependiendo de si el término está presente, o no, en el documento; en cuyo caso a la matriz se la denomina como \emph{matriz documento-término binaria}.
Una segunda opción de granulo más fino sería, en lugar de valorar a cada término simplemente por su ausencia o presencia en el documento, hacerlo por la cantidad de veces que se hizo presente. De esta forma, los términos más frecuentes en un documento serán ``más pesados'' (\emph{i.e.} más valiosos) que los menos frecuentes.
Sin embargo, al utilizar este pesado, el valor de los términos de cada documento estará ligado directamente a su tamaño.\footnote{Pensar, por ejemplo, el siguiente documento pequeño de cuatro palabras: ``la clasificación de texto''. Cada palabra tiene una frecuencia igual a 1, ahora pensar la frecuencia de estas mismas cuatro palabras, pero esta vez utilizando a este informe de tesis completo como documento.}
Por este motivo, raramente es utilizada la frecuencia en bruto como tal, no sin antes ser normalizada de alguna manera.
De hecho, en la práctica, se suelen utilizar numerosos métodos de normalización sobre estas frecuencias en bruto, tales como \emph{logaritmo de la frecuencia}, \emph{entropía}, \emph{ganancia de información}, \emph{divergencia de aleatoriedad}, \emph{BM25 score}, \emph{pesado Glasgow}, \emph{información mutua}, entre otros\cite{miner2012practical}.
Sin embargo, el método de normalización más utilizado es, sin lugar a dudas, el  \emph{tf-idf}, del inglés ``term frequency – inverse document frequency''\cite{salton1988term}, el cual dado un documento $d_i$ y un término $t_j$, calculará su peso $w_{i,j}$ como el producto de dos pesos especiales, $tf_{i,j}$ y $idf_{j}$, del siguiente modo:

\begin{equation}
w_{i,j}=tf_{i,j}\times idf_{j}
\label{eq:tf-idf}
\end{equation}

Donde $tf_{i,j}$ es la frecuencia normalizada del término $t_j$ en el documento $d_i$, esto es:\footnote{En la práctica esta normalización se puede llevar a cabo de distintas maneras, aquí sólo daremos la normalización por longitud del documento, que es una de las más utilizadas.}

$$tf_{i,j}= \frac{fr_{i,j}}{\sum^m_{k=1}fr_{i,k}}$$

Y donde $idf_{j}$ es (el logaritmo de) la inversa del número de documentos en los que el término $t_j$ aparece, en símbolos:

$$idf_{j}=\log\left(\frac{N}{df_j}\right)+1$$

En donde $df_j$ denota el número de documentos en los que el término $t_j$ aparece.

\

En la práctica, este método de pesado suele desempeñarse consistentemente bien, sin depender demasiado de la tarea que se aborde ---de allí su popularidad. Esto probablemente se deba a que su forma de valorar los términos, dada por la \autoref{eq:tf-idf}, parece encontrar un buen equilibro entre las siguientes dos intuiciones: cuanto más a menudo aparece un término en un documento, más representativo de su contenido es (capturado por $tf_{i,j}$) y, en cuantos más documentos aparece un término, menos discriminativo es (capturado por $idf_{j}$).

\


Antes de finalizar con la temática de las representaciones dadas por léxico, vale la pena señalar que, generalmente, los vocabularios de términos son considerablemente grandes, incluso podrían ser exponencialmente grandes usando \emph{n-gramas},\footnote{En relación al valor de $n$, \emph{i.e.} cuadrático para bigramas, cúbico para trigramas, y así.} en casos extremos.
Esto quiere decir que los vectores representando los documentos serán igualmente grandes ---por ejemplo, si hay 100.000 términos en el vocabulario, todos los vectores tendrán una longitud de 100.000.
Sin embargo, la mayoría de los documentos no contendrá, naturalmente, la mayoría de los términos ---\emph{i.e.} en otras palabras, la mayoría de los vectores estarán compuestos en su mayoría por ceros. 
Como consecuencia, en la práctica, las matrices documento-término son típicamente grandes y dispersas\footnote{Es decir, la mayor parte de sus elementos serán cero.} por naturaleza, lo que puede ser desfavorable, no sólo desde un punto de vista de costo y recursos computacionales, sino también desde el rendimiento del clasificador.\footnote{Puede que algunos modelos, dada su matemática, no sean óptimos para trabajar con vectores dispersos.}
Una práctica habitual utilizada para mitigar este problema, denominada como ``reducción de dimensionalidad'', consiste en la utilización de distintos enfoques para reducir el tamaño de los vectores.
Estos enfoques se dividen en dos grandes familias, ``selección de características'' y ``extracción de características'', dependiendo respectivamente de si los vectores se reducen al seleccionar un subconjunto de sus atributos (formado por los atributos más relevantes), o si, de lo contrario, se extrae una nueva representación más compacta desde la representación original ---\emph{i.e.} en otras palabras, si se utilizan métodos para crear un vector más pequeño, completamente nuevo, en función del original.
En la siguiente subsección se describe representaciones pertenecientes a este último grupo.


\subsection{Representaciones dadas por Contexto}

La extracción de atributos semánticos que vayan más allá de las unidades léxicas de un texto no estructurado probablemente sea uno de los problemas más desafiantes del \acf{PLN}.
Estos atributos pueden referir al significado, sentido, interpretación o coherencia de los elementos textuales de interés.
Para lograrlo, es necesario realizar un análisis profundo de estos elementos textuales y, si bien es un problema abierto en el \ac{PLN}, existen algunos enfoques interesantes que buscan modelar este tipo de información semántica.
La mayoría de los enfoques que existen en la literatura se basan en el principio de la hipótesis distributiva, el cual establece que \emph{palabras con significados similares ocurren en contextos similares}\cite{sahlgren2008distributional}.
En este sentido, \emph{la extracción de atributos semánticos se restringe a estrategias de representación que modelan el significado de las palabras en base a (las palabras de) su contexto}.
De forma tradicional, esto se lleva a cabo suponiendo que, en la tarea abordada, existe un número determinado de conceptos (o temas) latentes, y, luego, utilizando la información de ocurrencia y co-ocurrencia de las palabras, se infiere su contribución a cada uno de ellos\cite{lavelli2004distributional}.
Por ejemplo, el \emph{Análisis Semántico Latente}, más conocido como \emph{LSA} del inglés \emph{``Latent Semantic analysis''}, es uno de los métodos más conocidos para modelar la semántica de los documentos en relación a un conjunto cerrado de conceptos abstractos, típicamente denominados como ``temas latentes'' o simplemente ``temas'' \cite{deerwester1990indexing, dumais2004latent}.
Con LSA, una vez fijado un número, $r$, de temas latentes de interés,\footnote{En otras palabras, una vez elegida la longitud deseada para los vectores representando los documentos.} primero se aplica una \emph{descomposición en valores singulares} truncada\footnote{Truncada a $r$ valores singulares, es decir, tras la factorización de una matriz $M$ de $n \times m$, la matriz $\Sigma$ conteniendo los valores singulares deberá ser de $r \times r$, en lugar de $m \times m$.} sobre la matriz documento-término, $M$, factorizándola del siguiente modo:

$$M=U\Sigma V^T$$

Luego, la representación de los documentos estará dada por la matriz resultante del siguiente producto:

$$M'=U\Sigma$$

En palabras, primero se descompone la matriz documento-término original en tres matrices, $U$, $\Sigma$ y $V$, cada una de las cuales captura cierta información no explícita (latente) en la matriz original $M$ ---en concreto, puede verse a $U$ como capturando la distribución de los ``temas latentes'' en relación a los documentos, a $\Sigma$ como la importancia de cada tema y a $V$ como la pertenencia de cada término a cada tema.
Luego, al aplicar el producto de las dos primeras, $U$ y $\Sigma$, obtenemos una nueva matriz $M'$, la cual ya no es una matriz documento-término como $M$, sino una matriz ``documento-\emph{tema}''. Esta nueva matriz $M'$, al ser producto de $U$ y $\Sigma$, puede interpretarse como conteniendo el grado de pertenencia de cada documento a cada ``tema latente'' (heredado de $U$), normalizado por la importancia relativa de cada tema (heredada de $\Sigma$).
De este modo, utilizando esta nueva matriz, cada documento es representado como un vector de $r$ dimensiones ``semánticas'' (temas latentes), donde $r$ es menor que $m$, el número de términos en el vocabulario.

\

Finalmente, es importante mencionar que, con el auge del aprendizaje de representaciones y el \emph{Aprendizaje Profundo}, actualmente hay una tendencia cada vez mayor hacia obtener este tipo de representaciones semánticas aprendiéndolas directamente desde los datos, y no aplicando transformaciones sobre la matriz documento-término.
Concretamente, se utilizan técnicas y modelos tales como \emph{Universal Sentence Encoder}\cite{cer2018universal}, \emph{BERT}\cite{devlin2018bert}, \emph{Word2Vec}\cite{goldberg2014word2vec}, \emph{Doc2Vec}\cite{le2014distributed}, \emph{GloVe}\cite{pennington2014glove}, y \emph{FastText}\cite{joulin2016bag}, los cuales utilizan distintas arquitecturas de Redes Neuronales para aprender representaciones que, bajo el mismo principio de la hipótesis distributiva, intentan capturar la semántica en base al contexto en el que aparecen las palabras.
El uso del término \textbf{embedding} para referirse tanto a estas representaciones como a los vectores resultantes de las mismas se ha vuelto una norma.

% -Una vez elegidos los atributos, podemos aplicar cualquier modelo de clasificación que hemos visto; los populares para categorización de texto incluyen KNN, SVM, DT, MNB, y LG. En la siguiente sección veremos un escenario especial de clasificación de texto, en el que ...

\section{Clasificación Anticipada de Texto}
% \section{Analysis of Sequential Data: Early classification}
\label{sec:EC}

Hasta aquí se han introducido, y examinado, numerosos conceptos que sirven como el marco teórico general en el que el presente trabajo de investigación está ubicado.
En esta última sección, utilizando dicho marco, se examinará específicamente el escenario de clasificación de texto abordado en este trabajo, el cual, en su forma más general, se ubica dentro del área de análisis de datos secuenciales. %, denominado como \emph{clasificación anticipada de texto}.
Esta área de investigación es muy activa y aborda problemas en los que los datos son procesados naturalmente como secuencias, o pueden ser mejor modelados de esta manera, tales como la traducción automática~\citep{cho2014learning}, el análisis de videos y secuencias de imágenes~\citep{ullah2018action,srivastava2015unsupervised}, el reconocimiento de voz~\citep{graves2013speech} o el procesamiento de series de tiempo~\citep{kadous2002temporal}.
% Esta área de investigación es muy activa y aborda problemas en los que los datos son procesados naturalmente como secuencias, o pueden ser mejor modelados de esta manera, tales como el análisis de sentimientos~\citep{tang2015document}, la traducción automática~\citep{cho2014learning}, el análisis de videos y secuencias de imágenes~\citep{ullah2018action,srivastava2015unsupervised}, el reconocimiento de voz~\citep{graves2013speech} o el procesamiento de series de tiempo~\citep{kadous2002temporal}.
Los modelos y técnicas que se han utilizado para el análisis de datos secuenciales han incluido la técnica \emph{Dynamic Time Warping}~\citep{muller2007dynamic,sabinas2013one}, modelos ocultos de Márkov~\citep{yi2009hidden}, Procesos de decisión de Markov parcialmente observables~\citep{dulac2011text,rafferty2016faster} y diferentes tipos de \ac{RNR} como las redes \emph{\ac{LSTM}}~\citep{hochreiter1997long}, LSTMs bidireccionales~\citep{graves2005framewise} o, también, redes \emph{Gated Recurrent Units(GRU)}~\citep{chung2014empirical}.


Aunque estas secuencias pueden ser analizadas de diferentes maneras~\citep{gravessupervised}, en este trabajo nos centraremos en la \emph{clasificación} de las mismas~\citep{xing2010brief}. %, y en partícular, en una extensión específica de la misma, la \emph{clasificación anticipada}~\citep{xing2010brief}.
% A diferencia de la tarea de clasificación en vectores de atributos, las secuencias no tienen atributos explícitos. Incluso con técnicas sofisticadas de selección de características, la dimensionalidad de los potenciales atributos puede seguir siendo muy alta y la naturaleza secuencial de los atributos es difícil de capturar.
En este contexto, a diferencia de la clasificación convencional como la que hemos visto hasta ahora, en donde la entrada es procesada como un todo y representada por un vector, aquí la entrada debe ser procesada secuencialmente, elemento a elemento.
Esto provoca que la clasificación de secuencias sea una tarea más desafiante que la clasificación convencional sobre vectores de atributos~\citep{xing2010brief}, ya que los modelos tienen que \textbf{poder trabajar incrementalmente}, teniendo la capacidad de clasificar la secuencia, potencialmente, en cualquier instante de tiempo.
Así, una extensión de la clasificación de secuencias convencional es la denominada como ``clasificación anticipada''\footnote{En inglés, ``early classification''.}~\citep{xing2010brief}, en la que el problema es clasificar la secuencia \emph{lo antes posible}, sin tener una pérdida significativa de rendimiento.
Por ejemplo, algunos trabajos han intentado abordar la clasificación anticipada de texto utilizando distintas técnicas, tales como clasificadores \ac{MNB} \cite{escalante2015early}, representaciones de documentos basadas en perfiles\cite{escalante2017early} o representaciones específicas como \emph{Multi-Resolution Concept Representation}\cite{lopez2018early}.
% solo se centran en hacer predicciones basadas en información parcial, pero no abordan cómo seleccionar el prefijo más corto para proporcionar una predicción confiable.
Sin embargo, estos enfoques se han centrado en medir el rendimiento de los clasificadores al usar sólo información parcial de los documentos, pero no abordan cómo seleccionar el prefijo más corto que proporcione una predicción confiable. Esto es, consideran qué tan bien se desempeñan los clasificadores cuando sólo se les proporcionan porcentajes incrementales de los documentos, pero no exploran ningún mecanismo para decidir \emph{cuándo} (qué tan pronto) la información parcial leída hasta el momento es suficiente para clasificar la entrada con cierta seguridad.
Notar que esto no es un punto menor debido a que, por ejemplo, en escenarios realistas y dinámicos, como ser la clasificación de usuarios en redes sociales, no es posible utilizar un porcentaje fijo del contenido del usuario para clasificarlo, ya que dicho contenido es dinámico y virtualmente infinito, creciendo en el tiempo conforme el usuario lo genera; en escenarios de clasificación anticipada como éstos, es esencial que los modelos sean capaces de abordar la pregunta: ``¿es el contenido generado por el usuario hasta el momento suficiente para clasificarlo con ciertas garantías?''.
De hecho, la clasificación anticipada puede considerarse como un problema concreto de objetivos múltiples, en el que el desafío es encontrar un balance adecuado entre ambos, la anticipación y la efectividad de la clasificación~\cite{xing2010brief}.
Por lo que, en tareas de clasificación anticipada, los modelos deben ser capaces de soportar, o facilitar, algún \textbf{mecanismo de toma de decisión anticipada}.

Las razones detrás de este requerimiento de ``anticipación'' pueden ser diversas, podría ser necesario porque la longitud de la secuencia es virtualmente infinita (\emph{e.g.} contenido generado por usuarios) o, por ejemplo, si se pueden obtener ahorros o beneficios de algún tipo (\emph{e.g.} ahorro de recursos computacionales) clasificando la entrada de manera temprana.
Sin embargo, consideramos que los escenarios más importantes y desafiantes ocurren cuando la demora en la decisión puede, además, implicar consecuencias  riesgosas.
Este escenario, abordado en este trabajo de investigación y conocido como \emph{``detección anticipada de riesgos''}, ha ganado un interés creciente en los últimos años con posibles aplicaciones en la detección de rumores~\cite{ma2015detect,ma2016detecting,kwon2017rumor}, detección de depredadores sexuales e identificación de texto agresivo~\cite{escalante2017early}, detección de depresión~\cite{losada2017erisk, losada2016test} o detección de terrorismo~\cite{iskandar2017terrorism}.
% The reasons behind this requirement of ``earliness'' could be diverse. It could be necessary because the sequence length is not known in advance (e.g. online scenarios as suggested above) or, for example, if savings of some sort (e.g. computational savings) can be obtained by classifying the input in an early fashion.
% However, the most important (and interesting) cases are when the delay in that decision could also have negative or risky implications.
% This scenario, known as ``early risk detection'' have gained increasing interest in recent years with potential applications in rumor detection \cite{ma2015detect,ma2016detecting,kwon2017rumor}, sexual predator detection and aggressive text identification \cite{escalante2017early}, depression detection \cite{losada2017erisk, losada2016test} or terrorism detection \cite{iskandar2017terrorism}.
El objetivo central que impulsa el reconocimiento anticipado de estas situaciones de riesgo es intentar garantizar que las o la persona afectada reciba la atención y el apoyo social que necesite, a tiempo, ya que de no ser así, en casos extremos, la demora podría costarle hasta la vida.
Al igual que con cualquier otra aplicación crítica en salud, finanzas o seguridad nacional, este es un dominio en el que los modelos no sólo deben desempeñarse adecuadamente, sino que, además, deben ser confiables.
Esto es, del mismo modo que un médico no operaría, en efecto, a un paciente simplemente porque ``un modelo dijo que sí'', en estos escenarios de riesgo se requiere de un grado de confianza considerable en los modelos.
Una de las herramientas más útiles e importantes a la hora de decidir si un modelo es confiable o no, es que los expertos humanos puedan comprender las razones detrás de sus predicciones ---ya que ellos tienen una ``intuición adquirida'' propia del conocimiento del dominio que es difícil de capturar en un modelo con simples medidas de evaluación.
Así, en este tipo de escenarios, los modelos deberían, idealmente, \textbf{ser capaces de justificar o explicar sus decisiones}, de algún modo.
% Sin embargo, la mayoría de los modelos de clasificación funcionan como cajas negras, tomando una representación de la entrada, en forma de vector, y produciendo una respuesta en base a procesarla como un todo atómico e indivisible.

\

Como se mencionó anteriormente, uno de los puntos desafiantes en la clasificación anticipada es que los modelos aprendidos, generalmente, no brindan orientación sobre cómo decidir el momento correcto para clasificar la secuencia con ciertas garantías de precisión, con la información leída hasta ese momento. 
De hecho, hasta donde sabemos, el enfoque presentado por Dulac-Arnold \emph{et al.} (2015) \cite{dulac2011text} es el primero en abordar una tarea de clasificación (secuencial) de texto como un \emph{\ac{PDM}} con virtualmente tres acciones posibles: leer (la siguiente oración), clasificar\footnote{En la práctica, esta acción es un conjunto de acciones, una por cada categoría \emph{c}.} y detenerse.
La implementación de este modelo se basa en el uso de clasificadores \emph{\ac{SVM}}, los cuales fueron entrenados para clasificar cada acción posible como ``buena'' o ``mala'' en base a un ``estado actual'', \emph{s}.
% The key issue in real early sequence classification is that learned models usually do not provide guidance about how to decide the correct moment to stop reading a stream and classify it with reasonable accuracy.
% As far as we know, the approach presented in \cite{dulac2011text} is the first to address a (sequential) text classification task as a Markov decision process (MDP) with virtually three possible actions: read (the next sentence), classify\footnote{In practice, this action is a collection of actions, one for each category \emph{c}.} and stop.
% The implementation of this model relied on using Support Vector Machines (SVMs) which were trained to classify each possible action as ``good'' or ``bad'' based on the current state, \emph{s}.
Este estado actual está representado por un vector, $\Phi(s)$, el cual almacena las representaciones \emph{tf-idf} tanto de la oración actual como de las oraciones previas, así como también la categoría que el modelo asignaría hasta ese momento.
Aunque el uso de \ac{PDM} es muy atractivo desde un punto de vista teórico, y será considerado como trabajo a futuro, el modelo propuesto no sería adecuado para tareas de detección de riesgo.
El uso de \ac{SVM} como clasificador en base al vector $\Phi(s)$ implica que el modelo sea una caja negra, ocultando no sólo las razones para clasificar la secuencia como tal, sino también las razones detrás de su decisión de detenerse anticipadamente.\footnote{Dado que esta decisión está impuesta en el modelo por una función de recompensa que depende, a su vez, de los distintos vectores $\Phi(s)$, por los distintos estados \emph{s}.}
Las mismas limitaciones pueden encontrarse en trabajos más recientes \cite{yu2017learning, yu2018fast, shen2017reasonet} donde también abordan el problema de clasificación anticipada como un problema de aprendizaje por refuerzo, pero, en su lugar, usando Redes Neuronales Recurrentes.
% This state was represented by a feature vector, $\Phi(s)$, holding information about the \emph{tf-idf} representations of the current and previous sentences, and the categories assigned so far.
% Although the use of MDP is very appealing from a theoretical point of view, and we will consider it for future work, the model they proposed would not be suitable for risk tasks.
% The use of SVMs along with $\Phi(s)$ implies that the model is a black box, not only hiding the reasons for classifying the input but also the reasons behind its decision to stop early\footnote{Since this is enforced by the reward function which in turn depends, for each state \emph{s}, on vector $\Phi(s)$.}.
% The same limitations could be found in more recent works \cite{yu2017learning, yu2018fast, shen2017reasonet} also addressing the early sequence classification problem as a reinforcement learning problem but using Recurrent Neural Networks (RNNs).

Finalmente, en Loyola \emph{et al.} (2018) \cite{loyola2018learning}, se considera la decisión de ``cuándo clasificar'' como un problema aparte a ser aprendido por sí mismo.
De esta forma, se entrenan dos clasificadores \ac{SVM}, uno para predecir las clases y otro para decidir cuándo dejar de leer el flujo de entrada.
Sin embargo, el uso de estos dos \ac{SVM}, nuevamente, dificulta comprender las razones detrás de, respectivamente, tanto la clasificación como la toma de decisión en relación a cuándo detenerse prematuramente.
Además, como veremos más en detalle en la \autoref{sec:analysisAndDiscussion} del siguiente capítulo, cuando se utilizan modelos convencionales para clasificar secuencias, tales como \ac{SVM}, la clasificación puede volverse ineficiente y no escalable a escenarios reales, debido a que se tiene que reconstruir o actualizar continuamente la representación (\emph{i.e.} el vector) de la entrada completa cada vez que se procesa un nuevo elemento de la secuencia.\footnote{Por ejemplo, ineficientes en tiempo si se tiene que reconstruir la representación desde cero cada vez que un usuario crea un nuevo contenido; o en espacio, si se tienen que mantener distintas matrices documento-termino de respaldo (\emph{e.g. con frecuencias de términos}) para poder actualizar incrementalmente la representación.} Esto se debe a que, como se explicó en el segundo párrafo de esta sección, la mayoría de los clasificadores convencionales no fueron pensados para \emph{trabajar incrementalmente}.
% Finally,~\cite{loyola2018learning} considers the decision of ``when to classify'' as a problem to be learned on its own and trains two SVMs, one to make category predictions and the other to decide when to stop reading the stream.
% Nonetheless, the use of these two SVMs, again, hides the reasons behind both, the classification and the decision to stop early.
% Additionally, as we will see in \autoref{sec:complexity}, when using SVM to classify a document, incrementally, the classification process becomes costly and not scalable, since the \emph{document-term matrix} has to be re-built from scratch every time new content is added.

\

En resumen, la detección anticipada de riesgos es una tarea muy desafiante, ya que los modelos deben poseer, en simultáneo, un conjunto de cualidades deseables para poder abordarla adecuadamente.
En primer lugar, como en toda tarea de análisis de datos secuenciales, a diferencia de la clasificación convencional, los clasificadores deben \emph{poder trabajar incrementalmente}, teniendo la capacidad de clasificar la secuencia, potencialmente, en cualquier instante de tiempo. En segundo lugar, al tratarse de una tarea específica de clasificación anticipada, los modelos deben ser \emph{capaces de soportar, o facilitar, algún mecanismo de toma de decisión anticipada}.
Finalmente, y en tercer lugar, al tratarse de una tarea de detección de riesgo, los modelos deben brindar ciertas garantías de confiabilidad, por lo que deben ser \emph{interpretables, capaces de explicar naturalmente las razones de sus decisiones}.
En el siguiente capítulo, en el contexto de una tarea de detección de riesgos específica, la detección anticipada de depresión en redes sociales, introduciremos un nuevo modelo de clasificación diseñado específicamente con el objetivo de integrar estas características naturalmente.
\hfill$\square$